\documentclass[11pt]{article}

\usepackage[margin=1in]{geometry}
\usepackage[T1]{fontenc}
\usepackage[utf8]{inputenc}
\usepackage{lmodern}
\usepackage{amsmath,amssymb}
\usepackage{float}
\usepackage{graphicx}
\usepackage{setspace}
\usepackage[hidelinks]{hyperref}

\hypersetup{
  pdftitle={A Geometric Substrate Model for Bounded Local Intelligence},
  pdfauthor={Christopher A. Tucker},
  pdfsubject={Event-driven embedded intelligence, geometric partitioning, and bounded local execution},
  pdfkeywords={geometric substrate, bounded local intelligence, autonomous data partitioning, embedded autonomy, event-driven computation, deterministic local control}
}

\raggedbottom

\title{A Geometric Substrate Model for\\
Bounded Local Intelligence}
\author{Christopher A. Tucker\\
\small \texttt{cartheur@pm.me}\\
\small \href{https://orcid.org/0000-0002-0172-7258}{ORCID: 0000-0002-0172-7258}}
\date{}

\newcommand{\paperfigure}[3]{%
  \begin{figure}[H]
    \centering
    \includegraphics[width=#2\linewidth]{#1}
    \caption{#3}
  \end{figure}
}

\begin{document}
\maketitle

\begin{abstract}
This paper presents a self-organizing geometric partitioning model as a
substrate concept for bounded local intelligence. The immediate object of
analysis is a geometric construction procedure, but the central contribution is
architectural: the reinterpretation of that finite, symmetric, and addressable
structure as a local execution and partitioning substrate for embedded systems
that must sense, decide, and act under constrained power, memory, or
connectivity conditions. The walker-constructor model describes a bounded
computational data-space whose construction history, adjacency, and interior
partitions can be interpreted as supports for deterministic local behavior,
retraceable addressability, and tightly bounded control. The paper does not
claim a complete contemporary implementation on any single hardware platform.
It instead argues that the original geometric model remains a credible
conceptual precursor for low-power embedded autonomy, event-driven local
execution, and massively parallel asynchronous substrate design.
\end{abstract}

\onehalfspacing

\section{Introduction}

Contemporary AI systems are often framed in terms of large centralized models,
cloud orchestration, and probabilistic computation at scale. That framing is
powerful, but it does not fully answer a different class of engineering
problem: how to organize bounded local intelligence in low-power embedded
systems that must continue sensing, deciding, and acting under constrained
connectivity or limited computational budgets. The relevant design features are
deterministic local behavior, tightly bounded control, event-driven execution,
persistent low-power operation, and graceful degradation at the edge. The
present paper is valuable in that context because it offers a geometric model
that can be read not only as an abstract partitioning idea, but as a
substrate-oriented description of finite, addressable, local computational
organization. In that sense, the paper remains centrally about a geometric
algorithm for autonomous data partitioning within self-organizing systems and
the broader theory of automata.

Self-organizing algorithms are an efficient means to order complex
systems~\cite{ashby1957,simula1999,guerin2004}. The iterative construction of a
geometric structure can therefore be treated as more than a drawing procedure.
It can also be treated as a method for producing an internally coherent
partitioned space whose symmetries, adjacency relations, and bounded regions
are known to the system that constructs it. That interpretation is especially
relevant when the aim is not unbounded general computation, but local
organization of memory, addressing, traversal, and control in a finite
embedded setting. On that reading, self-organizing motifs are not merely
interesting abstractions~\cite{pattee1995}; they can also serve as models of
how local intelligence retains structure and acts within it.

Geometry, its practical application, and the intellectual problems following
have varied little between the ancient Greek world and modernity~\cite{netz1999}.
With the advent of the computer, more computationally sophisticated algorithms
were proposed~\cite{graham1972}, and computational geometry became more
explicitly tied to complexity, representation, and machine procedure
itself~\cite{toussaint1992,shamos1975}. Yet simplicity remains attractive when
the practical objective is not maximal abstraction, but a compact and legible
construction that can be traversed, addressed, and extended with bounded cost.
This matters for edge and embedded intelligence because such systems often need
finite local order more than they need global computational generality.

This paper therefore revisits a geometric algorithm drawn using limited tools
of a straightedge and drafting compass by automata. The automata, in the form
of goal-seeking agents, obey a rule that dictates they draw the structure with
the minimum amount of traveling while visiting each vertex and edge at least
once but at most twice. The activity of the automata yields a self-organized
structure of equivalent and symmetric dimensions that can be leveraged as a
partitioned computational data-space of finite block size. In the present
revision, that structure is interpreted as a conceptual local substrate for
bounded addressability, partitioned execution, deterministic embedded
autonomy, and event-driven local control. The claim is intentionally bounded.
The paper does not argue that the original geometric algorithm already
constitutes a complete modern platform implementation. It argues instead that
the model provides a credible and still-useful substrate concept for the kind
of low-power local intelligence that current asynchronous embedded systems seek
to realize.

\section{Problem Context: Edge Sensing, Bounded Autonomy, and Local Data Partitioning}

The immediate engineering problem addressed by this paper is how to organize
data, control, and local execution inside a finite embedded system that must
continue operating under bounded power, memory, and communications conditions.
In such settings, the system often cannot assume persistent upstream
connectivity, cloud orchestration, or unlimited instruction-time expansion.
Instead, it must preserve a usable local ordering of information and action.
That requirement is familiar in cybernetics and self-organizing systems, where
the relation among state, regulation, memory, and action matters more than raw
computational scale~\cite{ashby1957,guerin2004}. It is also relevant to the
critique of strictly instruction-centered machine models when the operative
question is how a machine remains behaviorally coherent while processing
changing local conditions~\cite{vonneumann1966,hartenstein2008}.

The present paper approaches that problem through geometric partitioning. The
claim is not that geometry alone solves embedded intelligence. The claim is
that a geometric construction can define a finite local space with known
symmetries, traversable adjacency, and partitionable interior regions. Such a
space can then be interpreted as a substrate for local addressability and
bounded execution rather than only as a mathematical figure. That makes the
paper especially relevant to architectures that emphasize deterministic local
behavior, partitioned control, event-driven compute, and graceful operation
within constrained environments.

\section{Original Geometric Model}

The original model is a self-organizing geometric algorithm whose structure is
built by three goal-seeking agents. The agents are simple: they walk, mark, and
remember. Their task is to construct a connected object while minimizing
unnecessary traversal. The central rule is that each vertex and edge must be
visited at least once and at most twice. That rule matters because it makes the
construction finite, bounded, and explicit. The paper is therefore not merely
about a shape. It is about a procedure for generating a known finite topology
whose relations are available to the agents that construct it.

Consider Fig.~1, a closed spatial system consisting of $n$ vertices with $m$
edges, where $n = 10$ and $m = 21$. The shape is constructed by three
goal-seeking agents of the form $f : P^{*} \rightarrow A$. They possess a
drafting compass, pins, and elastic strings, and collectively produce a
triangular arrangement with a central reference point and symmetric outer and
inner paths.

\paperfigure{figures/fig1.png}{0.42}{Walker constructor drawn by agents.}

The construction begins at an arbitrary origin, denoted $(0,0)$, on a
two-dimensional Cartesian grid. One agent moves in the positive $y$ direction,
while the other two move in the negative $y$ direction and then diverge along
opposing $x$ directions. The first polygon is formed by pinning an elastic
string across vertices $(1,2,3)$, after which compass operations and rays are
used to generate the additional points, including the inner points $a$, $b$,
$c$ and the outer points 4, 5, and 6. A second polygon is then formed across
$(4,5,6)$, and the intersection of the principal rays defines the central point
$m$.

The procedural significance of this sequence is that the resulting form is not
only symmetric, but also traceable. The system knows the relations that gave
rise to the form because those relations are the form's construction history.
That is why the model is interesting for computation. The shape is not an
arbitrary diagram that must later be interpreted. It is a generated space whose
adjacency, centering, and partition logic are already implicit in the method of
construction.

\paperfigure{figures/fig2.png}{0.88}{Walker constructor progress within algorithm execution.}

The collective behavior of the agents exhibits a Hamiltonian-path-like concern
with efficient traversal, though the paper uses that comparison more as a
structural intuition than as a formal reduction result~\cite{angelini2012,skiena1998}.
What matters here is that the agents produce a bounded connected data-space in
less descriptive time than would be required to discover the same structure
blindly. This is where the model begins to turn from geometry toward substrate
logic.

\section{System Interpretation Through Feature-Based Embedded Substrate Logic}

The geometric model in this paper is best understood as a conceptual local
substrate: a finite addressable interior whose partitions, edges, and center
can stand in for known regions of memory, attention, or control. The relevant
features are specific rather than brand-defined: deterministic local behavior,
tightly bounded control, low-power persistence, event-driven execution, and
massively parallel asynchronous composition.

Several traits of the model make this interpretation plausible.

First, the structure is bounded. The algorithm does not generate an unbounded
search field; it generates a finite local space. That aligns with the
requirements of bounded local behavior rather than globally open-ended
computation.

Second, the structure is deterministic in construction. Given the same rules
and initial relations, the same ordered space is produced. That does not make
the system universally deterministic in every downstream application, but it
does make the substrate itself legible and stable.

Third, the structure is locally intelligible. The relations among vertices,
edges, and interior points are not opaque. They are part of the generated
organization. This is important for embedded systems that must make decisions
using locally available knowledge rather than remote global state.

Fourth, the structure is compatible with event-driven interpretation. The
agents do not need a massive centralized planner to create or use the form.
They act through bounded local rules and preserved relations. This is why the
model sits more naturally within an event-driven embedded architecture story
than within a cloud-first one.

The claim remains conceptual. The paper does not demonstrate a deployed
substrate implementation. It instead provides a geometric use-case for how such
a substrate can be imagined: as a finite partitioned local data-space in which
construction history, adjacency, and addressability can support bounded
intelligence. One possible implementation context for exploring these ideas is
a GA144-based experimental board such as `Volatco`, but the present paper is
not a hardware sales document and does not depend on any single board identity.

\section{Partitioning, Addressability, and Local Execution Semantics}

The computational value of the structure appears most clearly when the interior
regions are treated as partitions rather than as blank geometric areas. The
original paper already points in this direction by describing a partitioned
data-space and by assigning addresses in the form of vertex tuples plus local
block identifiers. This is the point at which the model begins to resemble a
local execution substrate.

\paperfigure{figures/fig3.png}{0.42}{Data partitions following walker algorithm.}

In the original description, non-integral regions bounded by rays and interior
edges can be treated as storage partitions. The example partition $\{5, 0, c\}$
contains a finite number of addressable blocks, and a local address may be
written in a form such as $\{5, 0, c : 37\}$. That addressing style matters for
three reasons.

First, it anchors data location in known geometric relations. The system need
not discover location abstractly after the fact if the partition is already
defined by its construction path.

Second, it supports bounded resolution. The partitions are finite, and their
internal granularity is adjustable by changing local construction parameters
such as $\theta$. This resembles a tunable local-resolution model rather than a
single rigid address map.

Third, it supports retraceable traversal. Since the vertices and edges are
known, the system can interpret access not merely as indexed lookup, but as
movement through a locally intelligible graph of relations. In a substrate
model, that offers a useful way to think about how sensing, control, and local
state might be organized without presuming a fully centralized global store.

For this feature-based embedded reading, the relevant execution semantics are therefore
modest but meaningful:

\begin{enumerate}
\item local regions can be treated as bounded state or data spaces
\item access can be interpreted through known adjacency and construction paths
\item resolution can be tuned without abandoning the overall local geometry
\item traversal cost is constrained by the finite graph of the structure
\end{enumerate}

These are not yet the semantics of a full instruction set or operating system.
They are the semantics of an architectural substrate concept.

\section{Scaling Toward Multi-Node Asynchronous Systems}

The current paper is not primarily a scaling paper. Even so, the model suggests
a path to scaling through continuity of architectural logic across multiple
cooperating systems rather than simple enlargement of a single monolithic
compute space.

The local geometric substrate described here can be treated as one bounded node
of intelligence. If multiple such bounded nodes were composed, the important
question would not only be how much more compute becomes available, but whether
the same local partitioning logic remains legible across cooperating modules.
That is the key distinction between this interpretation and a conventional
cluster analogy. The point is not merely ``more nodes.'' It is the preservation
of the same local architectural character while functions are distributed
across boards, modules, or embedded assemblies.

In that sense, the present paper contributes to the scaling story only
indirectly. It clarifies what a bounded local substrate might look like before
the system is distributed. A later distributed-systems treatment would need to
explain how multiple such local spaces coordinate state, communication, and
control. That larger composition problem is not solved here, but the conceptual
boundary between local substrate and distributed scaling is now explicit.

\section{Application Classes and Use Cases}

The cleanest use-case fit for this paper is not generalized AI, nor large-scale
model serving. It is bounded embedded intelligence in settings where local
structure, low-power persistence, and graceful behavior matter more than raw
global optimization.

The strongest application classes are:

\begin{enumerate}
\item compact field robots that require bounded local sensing and action
\item embedded supervisory controllers with known local operating limits
\item low-power sensing and control nodes that must persist under weak
connectivity
\item local safety, fallback, or anomaly-handling subsystems inside a larger
system
\end{enumerate}

What these applications share is that the value lies in the local device's own
behavior: how it keeps functioning, how it partitions local information, and
how it remains intelligible under constrained conditions.

The paper is weaker as a direct fit for infrastructure-scale distributed
systems, because the current model remains centered on one bounded structure.
Those larger settings belong more naturally to later work on modular
composition and inter-node coordination. The present manuscript should
therefore mention scaling as an extension path, but should not overstate it as
a solved result.

\section{Complexity, Constraints, and Implementation Notes}

The original paper's complexity discussion is suggestive but should be read
carefully. The adjacency interpretation does support the intuition that known
edge relations can be tested efficiently in matrix form and that graph
traversals over the resulting structure admit the usual $O(n + m)$ search
behavior when the graph is explicit~\cite{orlin2010}. It is also reasonable to
state that subdivision and dimensional extension increase structural cost. What
should be avoided is any implication that the paper has fully established a
complete modern performance model for a deployed substrate.

The most defensible bounded claims are these:

\begin{enumerate}
\item the structure is finite and explicitly constructed
\item adjacency and partition relations are preserved by construction
\item finite partitions admit retraceable local addressing
\item higher-dimensional or more heavily subdivided variants would increase
operational complexity
\item the model is better understood as an architectural substrate concept than
as a benchmarked runtime implementation
\end{enumerate}

Several implementation questions remain open. The manuscript does not specify a
complete memory model, arbitration scheme, synchronization discipline, or
hardware realization. It also does not yet define the exact boundary between
construction-time knowledge and execution-time knowledge. Those questions are
not fatal to the paper, but they should be acknowledged as part of the honest
scope of the argument.

\section{Comparison With Conventional Architectures}

The paper's strongest contrast is not with all mainstream AI, but with a more
specific computational assumption: that useful intelligence is best organized
only through centralized instruction streams or remote generalized compute
resources. Against that assumption, the present model proposes a different
emphasis. It begins with a finite local structure, a bounded set of partitions,
known adjacency relations, and a method by which the system can interpret its
own local space.

That does not mean conventional architectures are obsolete. It means they
optimize for different strengths. Cloud-first and highly centralized systems
optimize for scale, aggregation, and global coordination. The geometric
substrate proposed here optimizes for local legibility, bounded addressability,
and deterministic spatial organization. The comparison is therefore not one of
universal superiority. It is one of architectural fit.

This is also why the paper remains relevant in the age of AI. Even where modern
statistical models are powerful, there remain practical reasons to preserve
local autonomy, explicit state organization, event-driven compute, and
constrained behavior inside edge systems. The present paper contributes to that
discussion by offering a compact substrate metaphor that is inspectable,
finite, and locally organized.

\section{Conclusion}

This paper has reinterpreted an early geometric partitioning model as an
architectural use-case for bounded embedded intelligence. Its strongest fit is
as a substrate-oriented account of how a finite, self-organized, and symmetric
local space may support partitioning, addressability, and local execution
semantics.

The contribution is intentionally bounded. The paper does not claim that the
original walker-constructor algorithm already constitutes a complete modern
embedded platform. It does claim that the model remains useful because it
describes a finite local structure whose construction history, adjacency, and
partitions are knowable to the system that builds it. That makes it a credible
conceptual precursor for deterministic local behavior, bounded autonomy,
partitioned embedded control, and event-driven substrate design.

The paper is most convincing when read as a local substrate model. Its relation
to larger multi-node asynchronous systems is real but secondary: the current
manuscript suggests how bounded nodes might later scale by preserving
architectural continuity, but it does not yet provide a full distributed
systems account. The practical value of the paper therefore lies in clarifying
a still-relevant design intuition: that finite local intelligence can be
organized through explicit spatial relations rather than only through
centralized instruction or remote generalized compute. That intuition is the
main reason the model remains worth preserving, rewriting, and extending.

\begin{thebibliography}{99}

\bibitem{ashby1957}
W.~R. Ashby,
\textit{An Introduction to Cybernetics},
Chapman \& Hall, Ltd., London, 1957.

\bibitem{simula1999}
O.~Simula et al.,
``Analysis and modeling of complex systems using the self-organizing map,''
\textit{Neuro-Fuzzy Techniques for Intelligent Information Systems} (1999): 3--22.

\bibitem{guerin2004}
S.~Guerin and D.~Kunkle,
``Emergence of constraint in self-organizing systems,''
\textit{Nonlinear Dynamics, Psychology, and Life Sciences},
Vol.~8, No.~2, pp.~131--46, 2004.

\bibitem{pattee1995}
H.~H. Pattee,
``Artificial life needs a real epistemology,''
\textit{Lecture Notes in Computer Science, Proceedings of the Third European Conference on Advances in Artificial Life},
Vol.~929, pp.~23--38, 1995.

\bibitem{netz1999}
R.~Netz,
\textit{The Shaping of Deduction in Greek Mathematics: A Study in Cognitive History},
Cambridge University Press, Cambridge, 1999.

\bibitem{graham1972}
R.~L. Graham,
``An efficient algorithm for determining the convex hull of a planar set,''
\textit{Information Processing Letters},
Vol.~1, pp.~132--33, 1972.

\bibitem{toussaint1992}
G.~Toussaint,
``What is computational geometry?''
\textit{Proceedings IEEE},
Vol.~80, No.~9, pp.~1347--63, 1992.

\bibitem{shamos1975}
M.~I. Shamos,
``Geometric complexity,''
\textit{Proc. 7th ACM Symposium on the Theory of Computing},
pp.~224--33, 1975.

\bibitem{ruffa2002}
A.~Ruffa,
``The generalized method of exhaustion,''
\textit{International Journal of Mathematics and Mathematical Sciences},
Vol.~31, No.~6, pp.~345--51, 2002.

\bibitem{knuth1969}
D.~E. Knuth,
\textit{The Art of Computer Programming},
Vol.~2. Addison-Wesley, Upper Saddle River, New Jersey, 1969.

\bibitem{dorst2002}
L.~Dorst and S.~Mann,
``Geometric algebra: A computational framework for geometrical applications,''
\textit{IEEE Computer Graphics and Applications},
Vol.~22, No.~3, pp.~24--31, 2002.

\bibitem{vonneumann1966}
J.~Von Neumann,
\textit{The Theory of Self-Reproducing Automata},
University of Illinois Press, Urbana, 1966.

\bibitem{hartenstein2008}
R.~Hartenstein,
``The von Neumann syndrome and the CS education dilemma,''
\textit{The International Workshop on Applied Reconfigurable Computing (ARC 2008)},
Mar.~2008.

\bibitem{orlin2010}
J.~Orlin et al.,
``A faster algorithm for the single source shortest path problem with few distinct positive lengths,''
\textit{Journal of Discrete Algorithms},
Vol.~8, Iss.~2, pp.~189--98, Jun.~2010.

\bibitem{angelini2012}
P.~Angelini et al.,
``Testing the simultaneous embeddability of two graphs whose intersection is a biconnected or connected graph,''
\textit{Journal of Discrete Algorithms},
Vol.~14, pp.~150--72, Jul.~2012.

\bibitem{skiena1998}
S.~Skiena,
\textit{The Algorithm Design Manual},
Springer-Verlag, New York, 1998.

\bibitem{grandoni2006}
F.~Grandoni,
``A note on the complexity of minimum dominating set,''
\textit{Journal of Discrete Algorithms},
Vol.~4, Iss.~2, pp.~209--14, Jun.~2006.

\end{thebibliography}

\end{document}
